\section{Composing a numerical bridge across an adaptive calculation}
\label{exec17:s9-composition}

The decision to port one routine becomes more demanding when its output is
an input to the next routine. Consider a servicing analysis that transforms
observations, fits an outcome model, evaluates an offer contrast, and uses that
contrast to choose a subsequent diagnostic. Comparing only the final mean
discards the mechanism by which a small arithmetic difference can change the
diagnostic selected. Conversely, requiring identical machine bytes at every
stage can reject a useful implementation without showing any difference in
the decision. A proposal for a bridge needs a compositional statement between
these extremes. The following constructions are project derivations; their
exact finite references do not execute an upstream causal package
\citep{exec17_s9_composition}.

First fix the scientific input and a particular sequence of branches. Let the
source states satisfy $z_i=f_i(z_{i-1})$, for $i=1,\ldots,m$, and let the decoded
bridge states be $\widetilde z_i=\widetilde f_i(\widetilde z_{i-1})$. A state
includes the quantities needed by later stages, not merely the displayed
estimate. All norms and units below must be specified: a bound on a probability
cannot be inserted as if it were a bound on a monetary contrast. Suppose both
states remain in a declared domain $D_{i-1}$, the source map has Lipschitz
constant $L_i$ there, and the bridge has a uniform local discrepancy
$\|\widetilde f_i(z)-f_i(z)\|\leq\epsilon_i$ for $z\in D_{i-1}$.
If $\|\widetilde z_0-z_0\|\leq e_0$, the triangle inequality gives
\begin{equation}
 e_i=L_i e_{i-1}+\epsilon_i,\qquad
 \|\widetilde z_m-z_m\|\leq
 e_0\prod_{i=1}^{m}L_i+
 \sum_{i=1}^{m}\epsilon_i\prod_{j=i+1}^{m}L_j.
 \label{exec17:s9-composed-error}
\end{equation}
The empty product is one. The proof proceeds by comparing
$\widetilde f_i(\widetilde z_{i-1})$ first with
$f_i(\widetilde z_{i-1})$, then with $f_i(z_{i-1})$. Induction proves the
recursion and its expanded form. The domain condition is substantive. A source
trajectory inside a well-conditioned region does not show that its perturbed
trajectory stays there. The implementation needs a containing domain or an
explicit failure when the next enclosure leaves it.

For a transparent arithmetic example, take $z_0=0.12$,
$f_1(z)=2z+1$ and $f_2(z)=z/10$. The source result is $0.124$. Let input
discrepancy be at most $0.001$, stage-one error at most $0.002$, and
stage-two error at most $0.0001$. The composed bound is $0.0005$, attained
when these three errors have the positive sign. Simply adding the three
error allowances gives a different and less informative quantity because it
ignores amplification and contraction. These constants are stipulated for
the example. They are not certificates for a learner, a matrix inverse or
a floating-point library.

An adaptive search adds a second obligation. Suppose stage $i$ branches on
$g_i(z_{i-1})\geq c_i$, where the source comparison function has Lipschitz
constant $M_i$ and the bridge's direct comparison error is at most $\eta_i$.
The common branch is certified only if
\begin{equation}
 |g_i(z_{i-1})-c_i|>M_i e_{i-1}+\eta_i.
 \label{exec17:s9-composed-branch}
\end{equation}
This condition must be applied inductively on the actual common history.
After an uncertified branch, the two runs can ask different questions, so the
next same-function comparison may not exist. One cannot continue adding local
errors as if the branch sequence had remained fixed. The appropriate output is
a set of possible paths, a return to the original engine, or a declared lack
of equivalence. Equality at the boundary also requires preserving the source's
strict or non-strict comparison convention.

This distinction matters for discovery because a conditional-independence
decision determines which later conditioning sets are inspected. It also
matters for model selection: an apparently negligible change in a diagnostic
can select a different learner with a different uncertainty construction.
For these cases a useful porting experiment retains the intermediate state
enclosures and branch margins. Comparing final graphs without these traces
does not explain why agreement occurred, while a differing graph does not by
itself identify an arithmetic defect. The source may itself be close to a
decision boundary. The reuse choice should therefore depend on a declared
profile of stable paths, not a promise that all data admit the same numerical
certificate.

\section{Conditional laws, hidden state and adaptive calls}
\label{exec17:s9-adaptive-kernels}

Randomized software requires a statement about the law of a sequence of calls.
Let $H_i=(X_1,\ldots,X_i)$ be the decoded history after call $i$, including
statuses, fitted-object identities, branch choices and any persistent random
state needed by the next call. On finite state spaces write
$K_i(x\mid h)$ and $\widetilde K_i(x\mid h)$ for the source and bridge
conditional distributions of the next result given the same history $h$.
If the initial histories have the same law and every reachable conditional
kernel agrees, multiplication of conditional probabilities gives the same
joint path law. Equality of one-call marginals is insufficient because it
does not determine these conditional probabilities.

For example, a source draws one fair bit $B$ and returns $(B,B)$ in two
related calls. A bridge starts a fresh worker for each call and returns two
independent fair bits. Each call has the same marginal distribution. The
source returns agreeing bits with probability one; the bridge does so with
probability one half. Direct enumeration gives joint total variation distance
$1/2$. The missing state is the shared bit. The analogy in a bank analysis is
a fitted control model or common bootstrap draw that must be reused across
offer contrasts. Assigning a seed to each request is not enough if the
definition of a request has silently split a shared computation.

For approximate kernels define
$\delta_i=\sup_h \frac12\sum_x|K_i(x\mid h)-\widetilde K_i(x\mid h)|$,
where the supremum covers all common histories reachable while comparing the
two implementations. On a given common history, the distributions can be
paired so that they return the same value with probability
$\sum_x\min\{K_i(x\mid h),\widetilde K_i(x\mid h)\}$, at least
$1-\delta_i$. Put their common masses together, then distribute the residual
masses to complete this pairing. Repeating this construction until the first
disagreement yields
\begin{equation}
 d_{\rm TV}\{\mathcal L(H_m),\mathcal L(\widetilde H_m)\}
 \leq 1-\prod_{i=1}^{m}(1-\delta_i)
 \leq\sum_{i=1}^{m}\delta_i.
 \label{exec17:s9-path-tv}
\end{equation}
The first inequality follows because every event has the same indicator when
the coupled paths agree. The second follows by expanding the probability of
at least one disagreement. This is a sufficient bound, often conservative;
it does not assert that a measured discrepancy in a few requests bounds the
supremum over histories. If the initial laws differ, an additional initial
coupling factor is needed.

Adaptive stopping fits the same argument when the horizon is bounded in
advance. Add a stopped state that repeats itself through the remaining calls,
and include the stopping decision in the kernel. A cap on conditioning-set
size or on iterations then has an explicit role. With an unbounded horizon,
one additionally needs a summable error allocation or a controlled tail for
the number of calls. Stopping after apparently favorable numerical results
cannot be justified by applying a fixed-horizon bound only to completed paths.

The path bound and the arithmetic bound answer different questions. Tiny
shifts of continuous numerical outputs can have disjoint point masses and
total variation one even when their metric distance is tiny. For continuous
states, one can instead couple numerical discrepancies in a specified metric
and use the margin condition to certify a common discrete branch. To bound
the probability of a branch error, one must supply the probability mass within
the relevant margin and the probability that the arithmetic certificate fails.
A density or margin bound is an extra assumption, not a consequence of
observing a small error on one input. Thus the proposal should retain both the
numerical enclosure and the conditional law of the discrete choices, using
each only on the domain where its assumptions have been established.

\section{When retries select a different statistical result}
\label{exec17:s9-retry-selection}

A language bridge can preserve every successful arithmetic call and still
change the reported statistical law through its failure policy. Fix one
scientific request. Let $Y$ denote the result that an ideal computation with
random input would produce, with law $P$. Let $S$ indicate whether that attempt
completes and returns a usable result, and write
$q(y)=\Pr(S=1\mid Y=y)$. This representation does not assume a failed
attempt actually calculated or logged its latent $Y$. It describes the joint
law required to reason about selection. Assume $0<\bar q=\int q(y)P(dy)$.
Among successful attempts the reported law is
\begin{equation}
 P_{\rm success}(dy)=\frac{q(y)P(dy)}{\bar q},\qquad
 E(Y\mid S=1)-E(Y)=\frac{\operatorname{Cov}(Y,q(Y))}{\bar q},
 \label{exec17:s9-success-tilt}
\end{equation}
when the displayed expectations exist. Both identities follow by integrating
the joint success mass $q(y)P(dy)$. A completed-only comparison estimates the
first quantity. It estimates the intended source mean only under further
conditions such as constant completion probability, or zero covariance for
that particular mean. Zero covariance does not establish equality of whole
distributions.

Now allow at most $K$ attempts, redrawing both computational and operational
randomness independently each time under the same law. Return the first
successful result and otherwise return the explicit symbol $\bot$. Summing
over the position of the first success gives
\begin{align}
 \Pr(R\in dy)&=q(y)P(dy)\sum_{j=0}^{K-1}(1-\bar q)^j,
 \nonumber\\
 \Pr(R=\bot)&=(1-\bar q)^K,
 &E(N)&=\sum_{j=0}^{K-1}(1-\bar q)^j,
 \label{exec17:s9-redraw-law}
\end{align}
where $N$ counts attempted calls, including the final failed call when all
attempts fail. Conditional on a return, the result still has law
$P_{\rm success}$ for every $K$. More redraws improve completion but do not
remove outcome-dependent selection. Taking an unlimited number of attempts
can make completion approach one while retaining the same biased returned
law. If the same outage, memory pressure or data difficulty persists between
attempts, the independent-attempt premise itself fails.

Preserving computational randomness produces another procedure. Draw $Y$ once
and retry the operational execution of that same result. Suppose failures are
conditionally independent given $Y$, with the same $q(Y)$ on each attempt.
The chance that this result is eventually returned is
$q_K(y)=1-(1-q(y))^K$, and
\begin{align}
 \Pr(R\in dy)&=q_K(y)P(dy),\qquad
 \Pr(R=\bot)=\int(1-q(y))^K P(dy),
 \nonumber\\
 E(N)&=\sum_{j=0}^{K-1}\int(1-q(y))^jP(dy).
 \label{exec17:s9-frozen-law}
\end{align}
Completed results are weighted by $q_K$, not by $q$. If $q(y)>0$ almost
surely, this selection vanishes as $K$ grows without limit, although expected
cost can be unbounded. If $q(y)=0$ on a positive-mass region, those results
never appear. Retaining a seed cannot repair a deterministic lack of memory,
an unsupported data type or a nonterminating algorithm. Even conditional
independence can be inappropriate when retrying retains the same defective
process state.

Take a stipulated score error $Y\in\{-1,1\}$ with equal probabilities,
$q(-1)=1/4$, and $q(1)=1$. This is a selection example, not an empirical
claim about causal software. With two attempts, exhaustive enumeration of
signs and four equally likely operational states gives
\begin{center}
\begin{tabular}{lccc}
\hline
Procedure & Completion & $\Pr(Y=1\mid\text{return})$ & Returned mean\\
\hline
One attempt & $5/8$ & $4/5$ & $3/5$\\
Two independent redraws & $55/64$ & $4/5$ & $3/5$\\
Two retained-result attempts & $23/32$ & $16/23$ & $9/23$\\
\hline
\end{tabular}
\end{center}
Both two-attempt procedures have expected call count $11/8$, yet their
returned distributions and completion probabilities differ. A proposal that
specifies only the retry count and average cost leaves the statistical
procedure undetermined. If one score unit represents USD100 of net offer
value, the ideal mean is zero while completed-only means are USD60 and about
USD39.13. These figures demonstrate a possible reporting distortion; they are
not estimates of the value of any banking intervention
\citep{exec17_s9_composition}.

\section{Persistent difficulty and honest all-request comparisons}
\label{exec17:s9-persistent-failures}

The independent redraw law can be optimistic even when each attempt has the
same marginal success rate. Let $\Theta$ denote a persistent request or
machine regime. Conditional on $\Theta=\theta$, assume attempts are
independent with success probability $q_\theta$. For a randomly assigned
request the no-return probability is
\begin{equation}
 \Pr(R=\bot)=E\{(1-q_\Theta)^K\}
 \geq (1-Eq_\Theta)^K,\qquad K\geq2.
 \label{exec17:s9-persistent-retry}
\end{equation}
The inequality follows because $x\mapsto(1-x)^K$ is convex on $[0,1]$.
It compares a mixture of conditionally independent attempt sequences with
an artificial sequence that redraws its regime on every attempt. In the
extreme finite example, half the regimes always complete and half never do.
Every attempt has marginal completion probability one half, yet no number
of within-regime retries reduces the no-return mass below one half. The
independent-attempt formula would instead give $2^{-K}$.

For a banking pipeline, difficult histories may create persistent large
conditioning sets, rare treatment arms or repeated optimizer failures.
Changing the language runtime may change this difficulty-selection relation.
Therefore the all-assigned comparison should retain the fixed request set,
its completion indicators, intended retry budget, actual attempts, and the
meaning of randomness retained between attempts. A useful analysis separates
variation between requests from repeated operational attempts within one
request. A thousand retries on one easy request are not a thousand
independent tests of the admitted request population.

There is a direct nonidentification issue if failed results cannot be
reconstructed. Fix one latent target result per assigned request, as in the
retained-result procedure. Suppose a bounded target $Y\in[a,b]$, completion fraction
$r$, and completed mean $\mu_s$ are observed under one declared procedure.
Writing $\mu_f$ for the unobserved failed-result mean gives
\begin{equation}
 E(Y)=r\mu_s+(1-r)\mu_f
 \in[r\mu_s+(1-r)a,\ r\mu_s+(1-r)b].
 \label{exec17:s9-failure-bounds}
\end{equation}
Without additional restrictions, any value in this interval can be attained
by assigning an admissible mean to failed requests. This is an identification
bound, not a confidence interval: sampling error in $r$ and $\mu_s$ needs
its own treatment. If the target is unbounded, finite bounds do not follow
from completion accounting alone. A timeout cannot honestly be coded as
zero, and a successful-subset mean cannot be labeled all-request performance.
For the retained-result example, the interval is $[0,9/16]$, containing
the ideal mean zero at its lower endpoint. Applying this decomposition to
independent redraws while calling the target the first draw's result would
be invalid: a returned later draw is a different latent quantity.

These mean bounds concern a latent analysis result. An operating decision
also needs the actual fallback action when analysis fails. Let $V$ denote
the underlying payoff state, distinct from the randomized analysis result
$Y$. If $U(A,V)$ is
the value of the chosen action and $C$ the total computation cost, the
relevant service quantity is
\begin{equation}
 E\{S_K U(A_{\rm returned},V)+(1-S_K)U(A_{\rm fallback},V)-C\},
 \label{exec17:s9-service-value}
\end{equation}
with $S_K$ indicating completion under the specified retry scheme. Comparing
only returned scores omits fallback consequences and cost; treating every
failure as the worst possible value can also distort the comparison if an
effective existing policy remains available. A public-data prototype can
stipulate these values transparently and vary them in sensitivity analysis.
It cannot infer a bank's actual fallback value, outage frequency or lost
business from the stipulation.

\section{Correcting selection, and the price of the correction}
\label{exec17:s9-selection-correction}

When the completion law is known and positive, one can recover some intended
functionals without claiming that the bridge has reproduced the source law.
For retained-result retries let $q_K(y)>0$ and let $h$ be the functional of
interest. Count every originally assigned request and use the contribution
$W=S_K h(Y)/q_K(Y)$, assigning zero on failure without needing its unobserved
$Y$. Conditional expectation gives
\begin{equation}
 E(W)=E\{h(Y)\},\qquad
 \operatorname{Var}(W)=E\{h(Y)^2/q_K(Y)\}-[E\{h(Y)\}]^2,
 \label{exec17:s9-inverse-selection}
\end{equation}
when the second moment is finite. The denominator of the empirical average
is the number originally assigned, not the number completed. This identity
requires a callable or independently justified $q_K$ for returned results.
It cannot be estimated from completed records alone without further
information about failures.

In the two-attempt retained-result example, $q_K(-1)=7/16$ and $q_K(1)=1$.
The inverse-weighted mean has expectation zero, but its per-request variance
is $23/14$, compared with one for an ideal observed score. Thus selection
correction restores the mean under its premises while increasing the
information cost. If the negative-score completion probability is mistakenly
entered as $1/2$ instead of $7/16$, the corrected expectation becomes $1/16$.
This residual bias is calculable even though every successful numeric score
is exact. It is not a floating-point error.

More generally, replacing $q_K$ by $\widehat q_K$ creates bias
$E[h(Y)\{q_K(Y)/\widehat q_K(Y)-1\}]$. If $|h|\leq H$,
$\widehat q_K\geq\varepsilon>0$ and
$|q_K-\widehat q_K|\leq d$ uniformly, its absolute value is at most
$Hd/\varepsilon$. A small absolute error in a rare completion probability
can therefore create a large error in the correction. Finite variance also
needs control of $h^2/q_K$, stronger than pointwise positivity. A practical
alternative is to narrow the admitted profile, make failed results observable
through a reference worker, or report bounds. Clipping inverse weights changes
the target unless its induced bias is separately bounded.

First-success redraws require another normalization. Given a returned result
with law proportional to $q(y)P(dy)$,
\begin{equation}
 E_P h(Y)=
 \frac{E_{\rm success}\{h(Y)/q(Y)\}}
      {E_{\rm success}\{1/q(Y)\}}.
 \label{exec17:s9-redraw-weighting}
\end{equation}
The numerator equals $E_P h/\bar q$ and the denominator equals $1/\bar q$,
which proves the identity. A finite-sample ratio of averages is generally
not unbiased, even with the correct $q$. Applying the retained-result
formula to independent-redraw outputs confuses two distinct sampling laws.
The retry semantics must consequently be part of the statistical procedure,
not an interchangeable configuration beneath it.

The resulting acquisition rule is specific. An original worker is attractive
when it permits retaining scientific state and exposing failures without
altering the procedure. A new bridge is attractive only when its proposed
benefit survives a comparison of numerical composition, conditional joint
laws, failure selection and fallback consequences. If completion is unrelated
to the target and branches remain certified, the obligations can simplify.
If these premises are unknown, a public-data experiment can test a bounded
profile and estimate diagnostics, while the proposal records the remaining
assumptions. No finite example here establishes universal port equivalence,
qualified dependent-data uncertainty, or a deployed banking service.
